\documentclass[11pt]{article}
\usepackage[T1]{fontenc}
\usepackage[utf8]{inputenc}
\usepackage{lmodern}
\usepackage[margin=1in]{geometry}
\usepackage{amsmath,amssymb,amsthm,mathtools}
\usepackage{booktabs,longtable,array}
\usepackage{microtype}
\usepackage{xurl}
\usepackage[colorlinks=true,linkcolor=blue,citecolor=blue,urlcolor=blue]{hyperref}
\newtheorem{theorem}{Theorem}[section]
\newtheorem{proposition}[theorem]{Proposition}
\newtheorem{lemma}[theorem]{Lemma}
\newtheorem{corollary}[theorem]{Corollary}
\theoremstyle{definition}
\newtheorem{definition}[theorem]{Definition}
\newtheorem{remark}[theorem]{Remark}
\newcommand{\R}{\mathbb R}
\newcommand{\C}{\mathbb C}
\newcommand{\dd}{\,d}
\newcommand{\ii}{\mathrm i}
\newcommand{\Rea}{\operatorname{Re}}
\newcommand{\Ima}{\operatorname{Im}}
\newcommand{\wind}{\operatorname{wind}}
\newcommand{\ind}{\operatorname{ind}}
\newcommand{\erf}{\operatorname{erf}}
\newcommand{\erfc}{\operatorname{erfc}}
\newcommand{\sgn}{\operatorname{sgn}}
\DeclareMathOperator{\Arg}{Arg}
\title{Heat Flow Contours, Collision Indices, and Gaussian Cap Energy\\
\large Initial research manuscript: results and open estimates}
\author{Drafted for Edward Baker\\
\small With substantial LLM assistance from gpt-6-astra (Codex), effort ultra}
\date{20 October 2026}
\begin{document}
\maketitle
\begin{abstract}
We study finite space--time contours around a hypothetical multiple real
zero of the theta heat family. With the orientation \(dt\wedge dx\), the
physical jet \(J=H+\ii H_x\) has local index
\(-\lfloor m/2\rfloor\) at a zero of spatial multiplicity \(m\). An ordinary
initial collision can be removed by an explicit quotient with a paid
apex patch; negative quotient winding would then force a later collision.
We prove the local index formula, a coupled complete-source remainder,
an obstruction for unsigned disks centered at the normalized initial
jet, and exact root-wall and nonzero-level contour formulas. Polynomial
and positive-source controls show why heat-flow geometry and local jets
alone do not force an additional collision.

We also derive localized Laguerre identities, energy balances, exact
Gaussian cap moments, and a sharp Gaussian deficit inequality. A negative
fixed-scale deficit would exclude the initial event directly; no such
theta deficit has yet been certified. One existing bounded interval
calculation proves positivity of \(H\) on
\([0.099999,0.100001]\times[1001,1002]\); we describe its complete error
analysis and distinguish it from a signed winding estimate.
The later-collision estimate, broad parameter coverage, and the desired
upper bound on the de Bruijn--Newman constant remain open.
\end{abstract}
\tableofcontents

\section{Problem, normalization, and status}
\subsection{The complete theta source}
For \(u\ge0\), set
\begin{align}
 \lambda_n&=\pi n^2,&
 \phi_n(u)&=(2\lambda_n^2e^{9u}-3\lambda_ne^{5u})
                      e^{-\lambda_ne^{4u}},&
 \Phi(u)&=\sum_{n\ge1}\phi_n(u), \label{eq:source}\\
 H(t,z)&=\int_0^\infty e^{tu^2}\Phi(u)\cos(zu)\dd u. \label{eq:H}
\end{align}
Each \(\phi_n(u)>0\), since
\(\lambda_ne^{4u}\ge\pi>3/2\). On compact complex \(t,z\) sets, the
source and each differentiated integrand have an integrable
superexponential majorant. Thus \(H\) is jointly analytic, is real and
even in real \(x\), and satisfies
\begin{equation}
                         H_t=-H_{xx}.                 \label{eq:heat}
\end{equation}
All its spatial derivatives are bounded on the whole real axis for a
fixed compact time interval. The sign in \eqref{eq:heat} matters:
\(u(s,y)=H(T-s,a+y)\) satisfies the ordinary forward heat equation
\(u_s=u_{yy}\). Throughout, ``later'' means increasing physical parameter
\(t\), hence the backward-heat direction.

Here is a normalization check independent of conventions for \(H\).
Let \(\Theta(v)=\sum_{n\in\mathbb Z}e^{-\pi n^2v}\) and
\(\psi(u)=\sum_{n\ge1}e^u e^{-\lambda_ne^{4u}}\).
The Gaussian Fourier transform followed by periodization gives
\(\Theta(v)=v^{-1/2}\Theta(1/v)\).
Consequently \(P(u)=e^u\Theta(e^{4u})=e^u+2\psi(u)\) is even,
\(\psi'(0)=-1/2\), and direct differentiation gives
\begin{equation}
             \Phi=(P''-P)/16=(\psi''-\psi)/8.          \label{eq:reflection}
\end{equation}
In particular \(\Phi\) extends evenly to \(\R\).
Splitting the Gaussian Mellin integral at \(1\), using theta inversion
on its lower part, and substituting \(v=e^{4u}\) gives, for
\(s=(1+\ii x)/2\),
\[
 \pi^{-s/2}\Gamma(s/2)\zeta(s)
  =\frac1{s(s-1)}+8\int_0^\infty\psi(u)\cos(xu)\dd u.
\]
The identity first follows in its convergence region and then by
analytic continuation. Twice integrating \(\psi''\) by parts proves
\begin{equation}
 8H(0,x)=\xi((1+\ii x)/2),\qquad
 \xi(s)=\tfrac12s(s-1)\pi^{-s/2}\Gamma(s/2)\zeta(s).   \label{eq:xi}
\end{equation}

\subsection{The threshold question}
The classical de Bruijn--Newman theory supplies a finite real constant
\(\Lambda\) such that all zeros of \(H(t,\cdot)\) are real exactly when
\(t\ge\Lambda\) \cite{debruijn,newman}. By \eqref{eq:xi}, the Riemann
hypothesis is equivalent to \(\Lambda\le0\). Rodgers and Tao proved
\(\Lambda\ge0\) \cite{rodgerstao}. We use these established facts as
background; we do not reprove the threshold theorem.

The contemplated contradiction begins with an actual multiple real zero
\((T,a)\) at a hypothetical positive threshold \(T=\Lambda\).
For the finite execution round the parameter envelope was
\[
 0<T\le\tfrac15,\qquad a\in[1001,1002],\qquad
 H(T,a)=H_x(T,a)=0,\qquad H_{xx}(T,a)\ne0.
\]
This is a conditional class, not an observed event or a coverage theorem.
No uniform lower bound on \(|H_{xx}(T,a)|\) is assumed.
An event at a threshold is an additional premise of the proposed route;
this manuscript does not prove that every hypothetical positive
threshold has an attained finite-height ordinary collision.

At any strictly later time \(t>\Lambda\), a multiple zero is impossible:
the local heat-polynomial analysis below gives nonreal zeros just before
such an event, still at times greater than \(\Lambda\).
Therefore a rigorous later-collision certificate would contradict the
threshold premise. It is the certificate, not this implication, that is
missing.

\subsection{What has and has not been obtained}
The analytical results below are derivations and internally reviewed
research results. The rational polynomial contour tests and the small
theta box are recorded computations with stated source and error
contracts. The reviews used separate agents of the same model; they are
not specialist review or independent numerical implementations.

The main positive estimating improvement is the removal of a mixed
absolute-error term from a complete coupled source remainder. The main
negative results concern precisely specified information classes:
unsigned disks, local prescribed jets, generic positive sources, and
unpaid or purely geometric contour changes. They do not exclude all
possible theta-specific signed estimates.
No negative theta quotient winding, negative Gaussian deficit at a
selected fixed scale, negative localized Laguerre average, or later theta
collision has been proved. The finite positive-value certificate excludes
one small box only.

\section{Physical degree and local collision index}
Write
\[
 J=H+\ii H_x,\qquad E=H^2+H_x^2.
\]
For a positively oriented boundary in the \((t,x)\) plane, with
orientation \(dt\wedge dx\), define
\begin{equation}
 W_J(\Omega)=\frac1{2\pi}\oint_{\partial\Omega}
 \frac{(H_xH_{xx}-HH_{xxx})\dd t+(HH_{xx}-H_x^2)\dd x}{E}.
                                                        \label{eq:winding}
\end{equation}
This is the integer winding of \(J\), provided \(J\ne0\) everywhere on
the boundary. A collision on a boundary, including an apex, makes this
physical winding undefined until the contour is changed or regularized.

\begin{theorem}[Local collision index]\label{thm:index}
Let \(H\) be a nontrivial real-analytic solution of \eqref{eq:heat}
near \(p=(T,a)\). If its spatial zero at \(p\) has exact order \(m\ge2\),
then \(p\) is an isolated zero of \(J\), and
\[
                         \ind_pJ=-\lfloor m/2\rfloor .
\]
\end{theorem}
\begin{proof}
Put \(h_j=\partial_x^jH(T,a)\), and define
\[
 P_m(s,z)=\sum_{k=0}^{\lfloor m/2\rfloor}
              \frac{m!(-s)^kz^{m-2k}}{k!(m-2k)!}.
\]
This polynomial satisfies the same heat equation. If
\(w(z)=e^{-z^2/4}\), comparison of coefficients in
\(w(z-2v)/w(z)=e^{zv-v^2}\) gives
\[
 P_m(1,z)w(z)=(-2)^m w^{(m)}(z).
\]
Integration by parts shows that \(P_m(1,\cdot)\) is orthogonal, against
the positive weight \(w\), to every polynomial of degree less than \(m\).
If it had fewer than \(m\) real sign-changing zeros, multiplication by
the product of the linear factors at those zeros would give a polynomial
of one constant nonzero sign whose weighted integral is zero. This is
impossible. Hence all \(m\) roots are simple and real.
Scaling gives the same conclusion for \(s>0\).
For \(s<0\), the coefficients after extracting \(z^{m\bmod2}\) are
strictly positive as a polynomial in \(z^2\), including its constant term.
There are therefore \(m\bmod2\) simple real roots.
At \(s=0\), \(P_m=z^m\). Thus
\((P_m,(P_m)_z)\) vanishes only at \((0,0)\).

Take \(K=[-1,1]\times[-b,b]\), \(b=2m\).
For \(|s|\le1\) and \(|z|\ge b\),
\[
 \frac{|P_m(s,z)-z^m|}{|z|^m}
 \le\sum_{k\ge1}(m^2/b^2)^k=\tfrac13 .
\]
In particular the two spatial walls have nonzero real part, and all
real roots on the time faces lie inside those walls.
For any real \(C^2\) function \(f\) with nonzero endpoints and \(N\)
simple interior zeros,
\begin{equation}
 \Delta\arg(f+\ii f')
 =\arctan\frac{f'(v)}{f(v)}
  -\arctan\frac{f'(u)}{f(u)}-N\pi.               \label{eq:phasezeros}
\end{equation}
Indeed \(f'/f=(x-r)^{-1}+O(1)\) at a simple zero \(r\);
the principal arctangent jumps upward by \(\pi\), whereas the lifted
argument is continuous. Subtracting \(\pi\) at each crossing proves the
formula. On the two spatial walls the arctangent endpoint differences
are the entire phase increments. The four endpoint terms cancel,
with the upper time face traversed left to right and the lower face
right to left. Thus the polynomial jet has degree
\((m\bmod2-m)/2=-\lfloor m/2\rfloor\).

It remains to justify the passage from the polynomial to \(H\), including
isolation. Its convergent Taylor series
\(H(T+\tau,a+y)=\sum a_{kj}\tau^ky^j\) satisfies
\[
 a_{kj}=\frac{(-1)^kh_{j+2k}}{k!j!}.
\]
The leading parabolic weight \(2k+j=m\) is \(cP_m\), where
\(c=h_m/m!\ne0\); all lower weights vanish.
Choose \(r_0>0\) with absolute convergence on a larger polydisc than
\(|\tau|\le r_0^2,\ |y|\le r_0b\), also for the differentiated series.
The finite tail sums
\[
 C_0=\frac1{|c|}\sum_{2k+j\ge m+1}|a_{kj}|r_0^{2k+j-m-1}b^j,\quad
 C_1=\frac1{|c|}\sum_{\substack{2k+j\ge m+1\\j\ge1}}
                 j|a_{kj}|r_0^{2k+j-m-1}b^{j-1}
\]
give, uniformly on \(K\) for \(0<r\le r_0\),
\[
 \left|\frac{H(T+r^2s,a+rz)}{cr^m}-P_m(s,z)\right|\le rC_0,\qquad
 \left|\frac{H_x(T+r^2s,a+rz)}{cr^{m-1}}-(P_m)_z(s,z)\right|
 \le rC_1.
\]
The polynomial jet has a positive minimum norm \(\delta_m\) on
\(\partial K\). For \(r(C_0+C_1)<\delta_m/2\), the straight boundary
homotopy is nonzero. Every nonzero sufficiently small \((\tau,y)\)
belongs to such a scaled boundary, by taking
\(r=\max(\sqrt{|\tau|},|y|/b)\). This proves isolation as well.
Both the source scaling and the target normalization preserve
orientation: their determinants are \(r^3>0\) and
\(1/(c^2r^{2m-1})>0\). The degrees agree.
\end{proof}

\begin{corollary}[Complete event count]\label{cor:count}
If a bounded piecewise smooth Jordan domain has \(J\ne0\) on its boundary,
then
\begin{equation}
 W_J(\Omega)=-\sum_{p\in\Omega:J(p)=0}\lfloor m(p)/2\rfloor .
                                                        \label{eq:count}
\end{equation}
\end{corollary}
\begin{proof}
Analyticity and the heat equation exclude infinite spatial order unless
the solution is identically zero. Theorem~\ref{thm:index} gives isolation;
compactness and a zero-free boundary collar give finiteness. The form in
\eqref{eq:winding} is locally \(d\arg J\), hence closed. Remove small
disjoint event neighborhoods and apply Stokes' theorem. Their positively
oriented indices sum to the outer winding.
\end{proof}

For the earlier nonreal-root assertion, apply the scaled analytic
expansion also on compact complex \(z\) sets. At \(s=-1\) the
polynomial \(P_m(-1,z)=\ii^mP_m(1,z/\ii)\) has simple imaginary roots,
with at least one nonreal pair for \(m\ge2\).
Small disjoint disks around those roots, uniform convergence, and
Rouch\'e's theorem give nonreal roots of \(H(T-r^2,\cdot)\).
This proves the assertion used in the threshold discussion.

If a domain contains an initial event of multiplicity \(m_0\), no other
event at or before its time, and later event charge
\(K_{\mathrm{later}}=\sum_{t>T}\lfloor m(p)/2\rfloor\), then
\[
 W_J=-q_0-K_{\mathrm{later}},\qquad q_0=\lfloor m_0/2\rfloor.
\]
A complete upper bound \(W_J<-q_0\) would force a later event.
An interval excluding every integer instead contradicts a proposed
winding enclosure itself; adding events does not repair such an interval.

\section{Ordinary-event regularization and the coupled source remainder}
Assume now \(h_2\ne0\), set \(\tau=t-T,\ y=x-a\), and define
\begin{equation}
 L(\tau,y)=-h_2\tau+\ii(h_2y-h_3\tau),\qquad Z=J/L.       \label{eq:L}
\end{equation}
The real matrix of \(L\) has determinant \(-h_2^2\), so \(L\) vanishes
only at the origin and has index \(-1\).

\begin{proposition}[Complete coupled Taylor bound]\label{prop:coupled}
Suppose \(H(T,a)=H_x(T,a)=0\) and \(h_2\ne0\).
For \(0\le\tau\le R\), define the complete finite moments
\[
 M_k(T,R)=\int_0^\infty e^{(T+R)u^2}\Phi(u)u^k\max(1,u)\dd u.
\]
Then
\begin{equation}
 |J(T+\tau,a+y)-L(\tau,y)|
 \le \tfrac12(M_4\tau^2+M_2y^2).                      \label{eq:coupled}
\end{equation}
Let \(C=\max(M_2,M_4)/2\),
\(c_L=h_2^2/\sqrt{2h_2^2+h_3^2}\), and
\(\rho=\sqrt{\tau^2+y^2}\). The quotient extends continuously by
\(Z(0,0)=1\), and
\[
 |Z-1|\le C\rho/c_L,\qquad
 \rho\le c_L/(2C)\ \Longrightarrow\ \Rea Z\ge\tfrac12.
\]
\end{proposition}
\begin{proof}
Use the single complex source kernel
\[
 K(\tau,y;u)=e^{\tau u^2}
       \{\cos((a+y)u)-\ii u\sin((a+y)u)\}.
\]
The event equations and heat equation imply that \(J-L\) is the integral
against \(e^{Tu^2}\Phi(u)\) of
\[
 K(\tau,y)-K(0,0)-\tau K_\tau(0,0)-yK_y(0,0).
\]
The cancellation is inside this complete integral.
Put \(A=\tau u^2,\ B=yu\). The two real coordinates of
\(\frac{d^2}{ds^2}K(s\tau,sy;u)\), before multiplying the second by
\(-u\), are \(e^{sA}\) times a rotation of
\((A^2-B^2,2AB)\). Their norm is
\(e^{sA}(A^2+B^2)\).
The factor \(\max(1,u)\) bounds the remaining diagonal map.
Taylor's integral remainder, integrated against \(1-s\) on \([0,1]\),
therefore proves \eqref{eq:coupled}.
The smallest singular value of \(L\)'s matrix is at least its determinant
magnitude divided by its Frobenius norm, namely \(c_L\).
Consequently \(|L|\ge c_L\rho\), and the claims follow.
\end{proof}

For \(T\le1/5,\ R=1\), the moments with weight \(e^{6u^2/5}\)
are uniform majorants. They are finite definitions, not numerically
acquired constants in the current record. Separate bounds on the three
second derivatives would give
\[
 e_{\mathrm{old}}=\tfrac12(M_4\tau^2+2M_3|\tau y|+M_2y^2).
\]
The exact improvement with the same moments is
\(e_{\mathrm{old}}-e_{\mathrm{new}}=M_3|\tau y|\).
This is an unsigned remainder improvement. It supplies no phase sign.
For a positive-source truncation, apply the same vanishing bracket to
the disjoint omitted-mode and omitted-spatial source regions.
Its tail payment is
\(\tfrac12(M_{4,\mathrm{tail}}\tau^2+
M_{2,\mathrm{tail}}y^2)\); a constant separate tail payment would lose
the needed apex vanishing.

The quotient need not be differentiable at the apex and need not solve
the heat equation. Continuity and the nonzero homotopy suffice.
The patch radius is event-dependent and cannot be made uniform without
additional curvature information.

\subsection{Two legitimate finite domains}
One physical domain, with \(e>0\), is
\[
 \Omega=\{(\tau,y):y^2/k-e<\tau<R\}.
\]
The initial event lies in its interior. Its upper cap runs left to right;
its parabolic boundary runs back from the right endpoint to the left.
Every cap, wall, endpoint and event must be counted.
Alternatively use the forward domain
\[
 D=\{(\tau,y):0<\tau<R,\ y^2<k\tau\}
\]
and the continuous quotient \(Z\). Remove a sufficiently small apex
neighborhood, prove the right-half-plane patch above, and glue it back.
Assume \(H\) is analytic near the closure and \(J\ne0\) on every
nonapex part of the boundary. Since the reference \(L\) has no zero away
from the apex, the complete regularized degree is
\[
                         W_Z(D)=-K_{\mathrm{later}}.   \label{eq:regularcount}
\]
Thus a negative quotient degree is precisely the additional-event target,
not a second count of the original collision.

\section{What local data and unsigned bounds cannot force}
\begin{proposition}[Unsigned-disk obstruction]\label{prop:disk}
If \(Z\ne0\) on a closed loop and \(|Z-1|\le1\) everywhere, then its
winding is zero. More generally, constraints consisting only of
pointwise disks centered at \(1\), together with continuity and
nonvanishing, cannot force negative winding.
\end{proposition}
\begin{proof}
The closed unit disk centered at \(1\) meets the imaginary axis only at
zero. Nonvanishing therefore puts the loop in the open right half-plane,
where an argument is single-valued. For disks of arbitrary nonnegative
radii, the constant loop \(Z=1\) satisfies every constraint and has
winding zero.
\end{proof}
The last statement is about that information class. It does not assert
that the constant loop is an admissible theta solution. Nor does it rule
out signed information correlated with those bounds.

There is also a heat-compatible control matching any prescribed
\(h_2\ne0,h_3\):
\begin{equation}
 H_c(\tau,y)=h_2(y^2/2-\tau)+h_3(y^3/6-\tau y).          \label{eq:cubic}
\end{equation}
If \(h_3=0\), this is quadratic, with threshold \(0\) and no later event.
If \(h_3\ne0\), set \(r=h_2/h_3,\ v=y+r\). Then
\[
 6H_c/h_3=v^3-3(r^2+2\tau)v+2r^3,
\]
whose discriminant is \(108((r^2+2\tau)^3-r^6)\).
Its sole real vanishing time is \(\tau=0\); it is positive for
\(\tau>0\). Thus the cubic has actual all-real threshold \(0\), the
specified ordinary collision, and no later collision.
Its exact remainder is
\[
 J_c-L=h_2y^2/2+h_3(y^3/6-\tau y)+\ii h_3y^2/2.
\]
Consequently the heat equation and these initial jets cannot force the
desired sign. This control lacks the complete theta source.

For the quadratic \(H=y^2-2\tau\), the maximum of the exact relative
remainder on the forward parabola and cap is
\[
 \max_{\partial D}|Z-1|
     =\frac{k\sqrt R}{2\sqrt{R+k}}.
\]
For example it exceeds \(1\) when
\(k>2,\ R>4k/(k^2-4)\), yet \(W_Z=0\).
Leaving the safe unit disk is not itself a winding mechanism.
Multiplying any of these fields by \(-1\) leaves winding unchanged.

\subsection{What a successful signed estimate must prove}
Let \(P=L+Q_{N,U}\) be a coupled finite-source proxy. A sufficient
certificate is
\[
 W(P/L)=-1,\qquad |J-P|\le B_{\mathrm{total}}<|P|
\]
on the entire nonapex boundary, with a separately paid compatible apex
homotopy. The winding and the strict separation both need certification.
Division does not create a separation advantage:
\[
 |Z-P/L|<|P/L|\quad\Longleftrightarrow\quad |J-P|<|P|.
\]
The comparison must use the same physical data and complete errors.

Explicit signed coordinates are
\[
 Z=(A+\ii B)/|L|^2,\quad
 A=-h_2\tau H+(h_2y-h_3\tau)H_x,\quad
 B=(h_3\tau-h_2y)H-h_2\tau H_x.
\]
At a simple negative-real-ray crossing \(B=0,A<0\), the contribution to
winding is \(-\sgn(dB/ds)\), for an oriented boundary parameter \(s\).
Every crossing, vertex and apex contribution must be controlled.
For the quadratic \(B=-2y^3\); its only nonapex candidates have
\(A=4\tau^2>0\), so no negative-ray crossing occurs.

\section{Finite contour controls and complete phase accounting}
\subsection{An exact rational consumer}
For a polynomially parametrized boundary cell, write each component of
the jet in its Bernstein basis on that cell. If all coefficients of one
component have a common strict sign, the whole cell image and its chord
lie in the same open half-plane. Replacing that image by its chord is a
nonzero homotopy. Subdivision yields a rational polygon with the same
winding.

Choose an integer shear \((X,Y)\mapsto(X,Y-qX)\) avoiding the real axis
at all polygon vertices. Such a \(q\) exists because only finitely many
values are forbidden, and the shear preserves orientation.
For an edge with endpoints of opposite \(Y\)-sign, the real-axis crossing
coordinate is
\[
 X_*=\frac{X_0Y_1-X_1Y_0}{Y_1-Y_0}.
\]
If \(X_*>0\), count \(\sgn(Y_1-Y_0)\); otherwise count zero.
The sum is the integer winding. Strict half-plane separation excludes an
edge through the origin. All operations in these polynomial tests are
rational. For nonpolynomial fields, an independently proved whole-cell
error bound is still required.

For a forward contour remove \(\tau<\eta\), with
\(\eta=b^2/k,\ b=1/32\). For the quadratic, \(|Z-1|\le\rho/2\),
so this removal has a nonzero quotient filling.
For the translated two-event quartic below, the local jet is
\(L=-12\tau+12\ii y\) and \(|J-L|\le59\rho^2\) for \(\rho\le1\).
At \(k=2,\ \eta=1/2048\), the removed patch has
\(\rho\le65/2048\), giving
\[
 |Z-1|\le3835/24576<1.
\]
Thus these quotient tests pay their removed apex rather than sampling
arbitrarily close to a physical boundary zero.

\subsection{Exact control fields}
Besides the quadratic, use
\[
 F_0(\tau,y)=y^4-12\tau y^2+12\tau^2.
\]
Its jet equations give either \(y=0\) or \(y^2=6\tau\); both force
\(\tau=0\). Its sole event is order four and has index \(-2\).
It cannot be passed to the ordinary quotient requiring \(h_2\ne0\).

The two-event control is
\[
 F_{-3}(s,y)=y^4-12sy^2+12s^2-3.
\]
It has exactly two ordinary events, at \((s,y)=(-1/2,0),(1/2,0)\).
The other critical branch \(y^2=6s\) gives
\(F_{-3}=-24s^2-3<0\).
Its squared roots are \(6s\pm\sqrt{24s^2+3}\); all roots are real
exactly for \(s\ge1/2\). Its first event is therefore not its threshold.
For the table translate the first event to \(\tau=0\), obtaining
\[
 H=y^4+6y^2-12\tau y^2+12\tau^2-12\tau .
\]
This is a detector control for an actual later event, not a counterexample
to threshold simplicity.

\begin{center}\small
\begin{tabular}{llllrr}
\toprule
Field & Boundary & \(k,R\) & Separation & Cells & Degree\\
\midrule
Quadratic & Physical & \(1,15/4\) & \(1/4\) & 7 & \(-1\)\\
Quadratic & Physical & \(2,7/4\) & \(1/2\) & 5 & \(-1\)\\
Quadratic & Physical & \(4,3/4\) & \(1/2\) & 6 & \(-1\)\\
Quadratic & Forward quotient & \(1,4\) & \(1/1024\) & 4 & \(0\)\\
Quadratic & Forward quotient & \(2,2\) & \(3/4096\) & 10 & \(0\)\\
Quadratic & Forward quotient & \(4,1\) & \(1/4096\) & 10 & \(0\)\\
Order four & Physical & \(4,3/4\) & \(3/4\) & 12 & \(-2\)\\
Two events & Physical & \(2,7/4\) & \(145/128\) & 12 & \(-2\)\\
Two events & Forward quotient & \(2,2\) & \(1/524288\) & 7 & \(-1\)\\
\bottomrule
\end{tabular}
\end{center}
For the physical contours \(e=1/4\) and the cap half-width is \(2\);
for the forward contours \(e=0\) and the same half-width is \(2\).
The stated separation is a certified physical-component half-plane
margin on the truncated cells, not a uniform lower bound on the
untruncated physical jet at its apex.
Each row was checked for both signs of the field, giving eighteen
successful cases. Three deliberate refusals were also checked:
a physical apex collision, a cap through the later event, and an
order-four event submitted to the ordinary quotient.
These tests validate the finite consumer on its controls; their number
does not provide evidence for the missing theta sign.

\subsection{Why parabolic wall phase cannot be dropped}
For \(H=y^2-2\tau\) on a physical contour, let
\[
 B^2=k(R+e),\quad A=B^2-2R,\quad
 \theta=\Arg(A+2\ii B)\in(0,\pi).
\]
The parabola and cap contribute respectively
\[
 I_{\mathrm{parabola}}=-2\theta,\qquad
 I_{\mathrm{cap}}=-2\pi+2\theta.
\]
Their sum is always \(-2\pi\). At fixed \(e>0\) and \(R\to\infty\),
the ordered pair tends to \((-2\pi,0)\) for \(k<2\),
\((-\pi,-\pi)\) for \(k=2\), and \((0,-2\pi)\) for \(k>2\).
For \(k<2\), the parabolic wall meets a simple root at
\(\tau=ke/(2-k)\). The jet remains nonzero there, but its wall phase
does not disappear. An infinite-cap argument needs a uniform bound on
every remaining boundary term and has not been supplied.

\section{Heat flow contours and collision-patch splitting}
\subsection{Diffusive scale without a causal cone}
An ordinary collision has born real roots
\begin{equation}
 x_\pm(T+\tau)=a\pm\sqrt{2\tau}
       +\frac{2h_3}{3h_2}\tau+O(\tau^{3/2}),\qquad \tau>0.
                                                        \label{eq:roots}
\end{equation}
Substitution into the Taylor expansion proves this formula; analytic
implicit continuation in \(\sqrt{\tau}\) justifies the expansions.
Before the event the leading pair is complex conjugate.
The square-root scale is diffusive, so ``heat flow contour'' is useful
terminology. Zero trajectories are not material paths of heat.
The Gaussian fundamental solution has nonzero tails at every distance:
there is no finite propagation cone \cite{nolen}.

For two subdomains separated by a nonzero internal cut, boundary
integrals add and the cut contributions cancel. A small rounded
neighborhood of one ordinary collision has winding \(-1\); a
collision-free remainder has total oriented boundary integral zero.
Such a remainder may be an annulus. Its outer loop can still have
winding \(-1\), balanced by the oppositely oriented inner loop.
Thus ``no collision in the remainder'' does not imply that every loop in
the remainder has zero phase circulation.

This distinction is visible in the exact quadratic. Take root walls
\(y=\pm\sqrt{2\tau}\), \(\eta\le\tau\le R\), a cap at \(R\), and a
lower connector
\[
 \tau=y^2-\eta,\qquad y:\sqrt{2\eta}\longrightarrow-\sqrt{2\eta}.
\]
The connector surrounds the apex and has
\(H=2\eta-y^2>0\) in its interior.
The walls have constant phase; the cap contributes \(-\pi\), and the
connector contributes \(-\pi\). The rounded lens has physical winding
\(-1\), exactly the original event. The remainder can be collision-free.
This realizes the proposed triangular split with no contradiction.

\subsection{Root-following walls}
\begin{lemma}[Cap phase between simple roots]\label{lem:rootcap}
Let \(f\) have simple endpoint roots and \(N\) simple interior roots on
\([u,v]\). Then
\[
                    \Delta_{u\to v}\arg(f+\ii f')=-(N+1)\pi.
\]
\end{lemma}
\begin{proof}
On each interval between consecutive roots \(f\) has constant strict
sign. If positive, the jet stays in the right half-plane, beginning
at positive imaginary axis and ending at negative imaginary axis.
Its argument change is \(-\pi\). If negative, multiply by \(-1\);
the increment is unchanged. Sum over all \(N+1\) intervals.
No pointwise monotonicity of the phase is needed.
\end{proof}
Along a simple root wall \(x=w(t)\),
\[
              w'=H_{xx}/H_x,\qquad J=\ii H_x.
\]
The phase is constant while simplicity persists.
Suppose the two born roots continue as compact simple walls to a finite
cap, and are adjacent at a sufficiently early lower cut.
If the upper cap has \(N\) interior simple roots, its phase is
\(-(N+1)\pi\), while the reversed lower cap contributes \(+\pi\).
Therefore
\[
 W_J(\hbox{truncated strip})=W_Z(\hbox{paid forward strip})=-N/2.
\]
The cap count \(N\) is even; adjoining the original event changes the
physical degree to \(-1-N/2\). To force a later event one needs \(N\ge2\).
This is additional information, not a consequence of following the roots.
Assuming continued adjacency would already set \(N=0\).
A failed tracker, escape from a chosen chart, or an unproved continuation
is not a certified collision.

For the two-event quartic, the outer walls are
\(\pm\sqrt{6s+\sqrt{24s^2+3}}\).
At \(s=1\) the inner roots are
\(\pm\sqrt{6-3\sqrt3}\), giving \(N=2\) and quotient degree \(-1\).

\subsection{Nonzero-level contours}
Let \(F=H/h_2\) near an ordinary event.
For \(e>0\) small, \(F=e\) has a local connected arc
\(\tau=y^2/2-e+\text{higher terms}\), since \(F_\tau(0,0)=-1\).
Suppose an actual such arc and an upper cap form a Jordan contour
enclosing the event, with exactly the intended two endpoints and
nonzero boundary jet. Along the arc
\[
                 J/h_2=e+\ii F_x,\qquad |J|\ge |h_2|e .
\]
Its argument change depends only on its endpoints.
If the left and right endpoint arguments in the right half-plane are
\(\theta_-,\theta_+\), and the cap has \(N\) simple interior zeros,
\eqref{eq:phasezeros} gives
\[
 I_{\mathrm{cap}}=\theta_+-\theta_- -N\pi,\qquad
 I_{\mathrm{arc}}=\theta_- -\theta_+.
\]
Hence the physical degree is \(-N/2\).
For one initial ordinary event and no other earlier event,
\[
                         N=2+2K_{\mathrm{later}}.
\]
The ordinary outcome is \(N=2\); the target requires \(N\ge4\).
Existence, continuation, cap endpoints and the event-dependent separation
\(|h_2|e\) all remain part of a usable numerical contract.
A graph \(\tau=\tau(y)\) can fail when \(H_{xx}=0\), without \(H=H_x=0\).
Failure of that chart is not an event.

More generally, writing \(Q=H_x^2-HH_{xx}\), the phase identities below
give constant-phase graphs \(x'=Q_x/Q\) where \(Q\ne0\).
Critical walls \(H_x=0,\ H\ne0\) also have constant phase.
Neither construction alone supplies a cap surplus.

\section{Other invariants, Laguerre expressions, and energy balances}
\subsection{Topological information already present in the jet}
The target of a nonzero jet is \(\C^*\), which retracts to the circle.
Its loop homotopy class is completely described by winding.
Passing to the unoriented real line of the jet gives the map
\(J^2/|J|^2\), of degree \(2W_J\), not new information.
Zero-crossing counts and root flux are equivalent descriptions under
their respective continuation hypotheses.
Embedding into a complex projective target can lose this invariant,
since \(\pi_1(\mathbb{CP}^1)=0\).
Spectral flow or an argument-principle contour for complex \(x\) may be
useful if supplied with a new operator or new complex-plane estimates;
they are not automatic refinements of this same boundary loop.

\subsection{Signed Laguerre form and localized averaging}
Direct differentiation gives, for any local argument \(\varphi=\arg J\),
\begin{equation}
 \varphi_x=-Q/E,\qquad
 \varphi_t=Q_x/E,\qquad Q_x=H_xH_{xx}-HH_{xxx}.          \label{eq:Qphase}
\end{equation}
For an even real entire function of order at most one with only real
zeros and nonzero value at zero, its paired canonical product is a
locally uniform limit of real-rooted polynomials.
For each such polynomial,
\(-(\log|H|)''\ge0\) away from zeros; by limits
\(Q\ge0\) everywhere.
The theta family has this growth bound: optimizing
\(r u+C u^2-c e^{4u}\) in its defining integral gives
\(\log M_H(r)=O(r\log(2+r)+\log^2(2+r))\).
Hadamard factorization, evenness and \(\sum r_j^{-2}<\infty\) then
give the paired product just described.
Thus a certified later point with \(Q<0\) contradicts later
all-realness directly. It does not by itself locate a later real
collision in a prescribed space--time domain.

Near an ordinary event the leading expression is
\[
 Q=h_2^2(\tau+y^2/2)+o(\tau+y^2),\qquad \tau\ge0.
\]
The analytic heat expansion makes the error uniform on a sufficiently
small forward parabolic neighborhood. Thus \(Q>0\) there away from the
event. A negative test must reach beyond this local region.
Moreover,
\[
                  Q_t+Q_{xx}=2(H_{xx}^2-H_xH_{xxx}),
\]
which introduces the next derivative form rather than a scalar
maximum-principle closure.

For a nonnegative compactly supported \(C^2\) weight \(\psi\),
integration by parts yields
\begin{equation}
 A_\psi:=\int\psi Q
       =2\int\psi H_x^2-\tfrac12\int\psi''H^2.          \label{eq:localized}
\end{equation}
On a finite interval the missing endpoint term is
\([-\psi HH_x+\frac12\psi'H^2]\).
When the integrals converge and \(HH_x\) vanishes at both infinite endpoints,
\(\int Q=2\int H_x^2\ge0\).
With \(w(u)=e^{tu^2}\Phi(u)\), symmetrization of products gives
\begin{equation}
 Q(t,x)=\frac14\iint w(u)w(v)
 \big[(u+v)^2\cos(x(u-v))+(u-v)^2\cos(x(u+v))\big]\dd u\dd v.
                                                        \label{eq:Qkernel}
\end{equation}
If the cosine transform of \(\psi\) is nonnegative at every frequency,
then \(A_\psi\ge0\) for every positive cosine source.
A centered Gaussian is such a weight. Translating the weight introduces
an oscillatory factor; it cannot be suppressed.
A fixed negative average would be a useful sufficient certificate.
Allowing an unspecified arbitrarily localized weight merely restates
the existence of a negative point.

\subsection{Energy identities and their limits}
Put \(A=H_x^2+H_{xx}^2\). The heat equation gives
\[
                    E_t+E_{xx}=2A.
\]
For moving endpoints \(a(t),b(t)\),
\[
 \frac{d}{dt}\int_a^bE
    =-[E_x]_a^b+2\int_a^bA+b'E(b)-a'E(a).
\]
The flux and moving-boundary terms are indispensable.
Along a simple critical branch \(H_x(t,b(t))=0,\ H_{xx}\ne0\),
\[
 b'=H_{xxx}/H_{xx},\quad
 V(t)=H(t,b(t)),\quad V'=-H_{xx},\quad
 (V^2/2)'=Q(t,b(t)).
\]
A fixed curvature sign controls the direction of \(V\) until the branch
degenerates; global continuation is an extra premise.

For an argument lift on a closed loop of length \(\ell\), Cauchy--Schwarz
gives
\[
 4\pi^2W^2\le\ell\int_0^\ell|\varphi'(s)|^2\dd s.
\]
This is an upper bound on \(|W|\), not a lower bound or a sign rule.
The zero-winding loop \(e^{\ii B\sin(2\pi s/\ell)}\) has arbitrarily
large angular energy. Conversely sufficiently small physical angular
energy could exclude an enclosed ordinary event; on a regularized loop
it only excludes additional charge.

A bulk form of the same degree is obtained from
\[
 \alpha_\varepsilon=
  \frac{H\dd H_x-H_x\dd H}{E+\varepsilon^2},\qquad
 d\alpha_\varepsilon=
 -\frac{2\varepsilon^2(H_{xx}^2-H_xH_{xxx})}
       {(E+\varepsilon^2)^2}\dd t\wedge\dd x .
\]
On a nonzero boundary Stokes' theorem and uniform boundary convergence
give
\[
 W_J=-\lim_{\varepsilon\downarrow0}\frac1\pi
 \int_\Omega
 \frac{\varepsilon^2(H_{xx}^2-H_xH_{xxx})}
      {(E+\varepsilon^2)^2}\dd t\dd x .
\]
This moves the same integer into a bulk integral. It neither justifies
a pointwise limit exchange in the interior nor repairs an unpaid apex.

\subsection{Finite zero-interaction energy}
For a polynomial heat solution with all roots real and simple,
differentiating \(H(t,x_i(t))=0\) and factoring the polynomial give
\[
 x_i'=H_{xx}(x_i)/H_x(x_i)
      =2\sum_{j\ne i}\frac1{x_i-x_j}.
\]
For the complete root system define
\[
 U=-\sum_{i<j}\log|x_i-x_j|.
\]
Then \(x'=-2\nabla U\) and
\[
                         U'=-2|\nabla U|^2\le0.
\]
In a bounded configuration a finite-time collision would make
\(U\to+\infty\), contradicting its decrease.
This explains forward repulsion after all roots are real.
For the infinite theta system one must specify renormalization,
exterior interactions and moving-cutoff terms. Applying the finite
formula only to a selected pair discards forces.
The utility of related energy methods in the lower-bound direction
\(\Lambda\ge0\) \cite{rodgerstao} does not provide the sought upper bound.

\section{Gaussian cap moments and a sharp energy test}
\subsection{Adjoint conservation and what heat transports}
If \(\psi_t=\psi_{xx}\), direct differentiation proves
\[
 \partial_t(\psi H)+\partial_x(\psi H_x-\psi_xH)=0.
\]
For simple moving root boundaries \(a(t),b(t)\), the signed lobe mass
\(M(t)=\int_a^bH(t,x)\dd x\) satisfies
\[
                  M'=H_x(a)-H_x(b).
\]
Its magnitude grows in increasing physical \(t\) for a single sign
lobe, consistently with the reversed heat direction.
For the local quadratic
\(H=h_2(y^2/2-\tau)\), it is
\(-4\sqrt2\,h_2\tau^{3/2}/3\).

For \(s>0\), set
\[
 G_s(y)=(4\pi s)^{-1/2}e^{-y^2/(4s)},\qquad
 f(y)=H(T+s,a+y),\qquad h_n=\partial_x^nH(T,a).
\]
Gaussian convolution multiplies \(\cos(uy)\) by \(e^{-su^2}\).
The superexponential source justifies Fubini and differentiation, so
\begin{equation}
 \int_\R G_sf=h_0,\qquad
 \int_\R G_sf'=h_1
     =\int_\R G_s\frac{y}{2s}f,\qquad
 \int_\R G_sf^{(n)}=h_n.                              \label{eq:capmoments}
\end{equation}
Thus a collision gives two exact vanishing moments of the later cap.
For a continuous nonzero \(f\), those two moments force at least two
sign changes: no sign change contradicts the mean-zero condition;
exactly one at \(b\) makes \((y-b)f(y)\) have a constant sign, contradicting
both moments. This is the initial pair's count, not an extra pair.

The quadratic root wedge \(|y|<\sqrt{2s}\) contains only
\(\erf(1/\sqrt2)\), about \(68.27\%\), of the Gaussian mass.
The omitted \(31.73\%\) does not shrink as \(s\downarrow0\).
For this quadratic the inner signed integral is
\(-2h_2s\,\phi(1)\), with \(\phi\) the standard normal density, and the
outer integral is its opposite. Truncating to the root triangle loses
exactly relevant heat data.
Level sets of the heat kernel instead give heat-ball sides
\(y^2=2s\log(A/s)\), \(0<s<A\), in the adjoint parameter \(s=t-T\).
They are natural mean-value geometries \cite{mccuan}, but are not
constant-phase walls and do not imply finite propagation.

\subsection{Hermite expansion and the deficit identity}
\begin{theorem}[Gaussian deficit]\label{thm:deficit}
For the complete theta family and any \(s>0\), define
\[
 D_s(T,a)=s\int_\R G_s f'^2-\int_\R G_s f^2 .
\]
Then the convergent identity
\begin{equation}
 D_s=-h_0^2-sh_1^2+
   \sum_{n\ge3}\frac{n-2}{2}\frac{(2s)^n}{n!}h_n^2       \label{eq:deficit}
\end{equation}
holds. At a common zero \(h_0=h_1=0\),
\begin{equation}
 \int G_sf^2\le s\int G_sf'^2\le2s^2\int G_sf''^2.       \label{eq:poincare}
\end{equation}
Equality in the first inequality at a nonzero event is possible only
for \(f(y)=c(y^2-2s)\).
\end{theorem}
\begin{proof}
Let \(z=y/\sqrt{2s}\), let \(\gamma\) be standard Gaussian measure,
and let \(\mathrm{He}_n\) be the monic Hermite polynomials defined by
\(e^{vz-v^2/2}=\sum \mathrm{He}_n(z)v^n/n!\).
Gaussian integration gives
\(\int\mathrm{He}_m\mathrm{He}_n\dd\gamma=n!\,\mathbf1_{m=n}\);
Rodrigues' formula gives the same identity by integration by parts.
These polynomials are complete in \(L^2(\gamma)\):
if \(g\) is orthogonal to all of them, then
\(\int g(z)e^{vz-v^2/2}\dd\gamma(z)\) is entire by
Cauchy--Schwarz, with all Taylor coefficients zero. Its values at
imaginary \(v\) show that the Fourier transform of \(g\gamma\) vanishes,
hence \(g=0\).

Expand \(f(\sqrt{2s}\,z)=\sum c_n\mathrm{He}_n(z)\).
Integration by parts and \eqref{eq:capmoments} give
\[
                    c_n=\frac{(2s)^{n/2}}{n!}h_n.
\]
The cap and all its derivatives are bounded, so belong to the needed
Gaussian Sobolev spaces. Parseval applied also to the derivatives gives
\[
 \int G_sf^2=\sum n!c_n^2,\quad
 s\int G_sf'^2=\tfrac12\sum n\,n!c_n^2,\quad
 2s^2\int G_sf''^2=\tfrac12\sum n(n-1)n!c_n^2 .
\]
The differences are convergent; subtraction proves
\eqref{eq:deficit}. At an event \(c_0=c_1=0\), and termwise comparison
proves \eqref{eq:poincare}. Equality in its first step eliminates every
coefficient except \(c_2\), giving the stated quadratic.
\end{proof}
The same proof applies to polynomial and finite positive atomic controls,
whose derivatives have adequate Gaussian integrability.
If \(h_0,\ldots,h_{m-1}\) vanish, it also gives
\[
               \int G_sf^2\le(2s/m)\int G_sf'^2.
\]
For the nonzero bounded theta cap, the quadratic equality case is
impossible; thus an event would have \(D_s>0\), without a uniform
positive margin being supplied.

A complete fixed-scale bound \(D_s\le-\delta<0\) would instead imply
\[
                         h_0^2+sh_1^2\ge\delta .
\]
It directly excludes the initial event, regardless of its multiplicity.
It is not a later-collision or negative-winding certificate.
If \(h_0\ne0\), \(D_s\to-h_0^2\) as \(s\downarrow0\).
If \(h_0=0,h_1\ne0\), the identity or a Taylor expansion gives
\(D_s=-sh_1^2+O(s^3)\).
Thus allowing an unspecified sufficiently small \(s\) merely
reformulates first-jet nonvanishing. A useful estimate must fix the scale
and parameter window before acquisition and be compared with a direct
consumer of the same data.

\subsection{Complete source kernel and exact controls}
Using
\(\int G_s\cos(y\omega)\dd y=e^{-s\omega^2}\) and product-to-sum
identities gives
\begin{align}
 D_s=\frac12\iint w(u)w(v)\big[&
 (suv-1)e^{-s(u-v)^2}\cos(a(u-v))\nonumber\\
 &-(suv+1)e^{-s(u+v)^2}\cos(a(u+v))\big]\dd u\dd v,
                                                        \label{eq:Dkernel}\\
 &w(u)=e^{(T+s)u^2}\Phi(u).\nonumber
\end{align}
Both terms, including their signs and correlated source data, must be
retained in a useful estimate.
For the quadratic control \(D_s=0\).
For \eqref{eq:cubic}, the only surviving term in \eqref{eq:deficit}
is \(n=3\), giving \(D_s=2h_3^2s^3/3\).

A bounded positive-source control is
\[
 H(t,x)=1+e^{t-T_0}\cos x,\qquad T_0>0.
\]
It has actual all-real threshold \(T_0\), with ordinary events at
\((T_0,(2k+1)\pi)\), and simple real zeros at all later times.
Indeed for \(\cos z=-e^{T_0-t}\), a nonzero imaginary part requires
\(\sin(\Rea z)=0\) and hence \(|\cos z|\ge1\);
this is impossible for \(t>T_0\). For \(t<T_0\) there are nonreal
solutions. At the event centered at \(\pi\), \(f=1-e^s\cos y\) and
\[
 D_s=1+s\sinh(2s)-\cosh(2s)>0\qquad(s>0).
\]
The expression and its first derivative vanish at zero; its second
derivative is \(4s\sinh(2s)>0\). This is a bounded positive cosine
source with no forced later event. Its simultaneous periodic events can
be isolated individually for local contours; it lacks theta's complete
arithmetic source. Positivity and boundedness alone therefore do not
supply the contrary deficit sign.

\subsection{Gaussian Bochner identity}
For \(F(s)=\int G_s(y)E(T+s,a+y)\dd y\), integration by parts and
\(E_t+E_{xx}=2A\) yield
\[
 F'(s)=2\int G_sA(T+s,a+y)\dd y,
\]
and hence
\[
 F(s)=E(T,a)+2\int_0^s\int G_rA(T+r,a+y)\dd y\dd r.
\]
At a quadratic event both sides after the initial term equal
\(2h_2^2(s+s^2)\). Subtracting the bulk term merely recovers the
initial nonnegative first-jet energy. It is an exact accounting identity,
not an independent sign contradiction.

\section{The bounded theta value certificate}\label{sec:numerical}
This section reports the completed bounded computation, rather than a
new acquisition. Its analytic formulas are included to specify exactly
what the recorded interval output proves.

\begin{proposition}[Recorded computational exclusion]\label{prop:box}
For the complete \(H\) in \eqref{eq:H}, the recorded outward interval
calculation gives
\[
 H(t,x)>4.1345\,10^{-4}G_*
 \quad\hbox{on}\quad
 \mathcal B=[99999/10^6,100001/10^6]\times[1001,1002],
 \qquad G_*=1000^{7/4}e^{-125\pi}>0 .
\]
In particular this closed box contains no zero of \(H\), and therefore no
multiple real zero of any multiplicity.
\end{proposition}
The computational premise is the validity of the outward source
quadratures and interval arithmetic described below. This is not a
symbolic derivation of the displayed decimal outputs.

\subsection{Complete contour shift before truncation}
Fix
\[
 d=\pi/8-4\log(1000)/1000,\quad m=12,\quad U=2,\quad
 t_c=21/200,\quad x_c=1000 .
\]
The mode series and reflection identity \eqref{eq:reflection} extend
normally to \(|\Ima u|<\pi/8\).
First write the complete even integral on \(\R\), shift it to
\(\R+\ii d\), and only then split and reflect the negative half.
Superexponential decay on the shifted strips makes the remote vertical
sides vanish. This gives
\begin{equation}
 H(t,z)=\frac{e^{-dz}}2\int_0^\infty
 \left[e^{t(u+\ii d)^2}\Phi(u+\ii d)e^{\ii zu}
       +e^{t(u-\ii d)^2}\Phi(u-\ii d)e^{-\ii zu}\right]\dd u .
                                                        \label{eq:shift}
\end{equation}
Individual theta modes are not even; their omitted vertical segments
cannot be discarded by asserting a finite-mode reflection.
Let \(F(t,z)\) denote \eqref{eq:shift} with \(1\le n<12\) and
\(0\le u\le2\). It is real on the real axis and satisfies the heat
equation. With \(v=u+\ii d\), its normalized midpoint derivatives are
\begin{equation}
 I_j=\frac{\partial_x^jF(t_c,x_c)}{G_*}
 =1000^{9/4}\Rea\int_0^2 e^{t_cv^2+v+\ii x_cu}
 \sum_{n=1}^{11}(2\lambda_n^2e^{8v}-3\lambda_ne^{4v})
 e^{-\lambda_ne^{4v}}(\ii v)^j\dd u .                   \label{eq:Ij}
\end{equation}
The inherited acquisition provides dyadic enclosures
\(I_j\in[a_j-\rho_j,a_j+\rho_j]\), \(0\le j\le50\).
These input enclosures are recorded in Appendix~\ref{app:cache}.

Let \(c=\cos(4d)\) and \(N_*=1000^{9/4}e^d\). A common normalized
majorant for derivatives through order \(52\) is
\begin{equation}
 \mathcal M=N_*\int_0^2\sum_{n=1}^{11}(2\lambda_n^2+3\lambda_n)
 e^{u^2/5+9u-c\lambda_ne^{4u}}[1+(u+d)^{52}]\dd u
 <9.632888121\,10^{10}.                               \label{eq:maj}
\end{equation}
This outward bound applies for \(0\le t\le1/5,\ x\ge999\):
\(N_*=e^{-999d}/G_*\), and the remaining spectral modulus
\(e^{-dx}\) decreases with \(x\). Thus the extension from an older
height edge at \(1001\) is analytically justified, not assumed.

\subsection{Taylor and source errors}
For \(j=0,1\), use
\[
 P_j(t,x)=\sum_{\ell=0}^{24}\sum_{k=0}^{12}
 \frac{(-1)^k a_{j+\ell+2k}}{\ell!k!}
          (x-x_c)^\ell(t-t_c)^k.
\]
These are component approximants; no identity
\(P_1=\partial_xP_0\) is assumed.
On a complete parameter cell put
\(R_x=\max|x-x_c|,\ \Delta=\max|t-t_c|\).
Spatial Taylor expansion followed by time expansion using the heat
equation gives
\begin{align}
 E_{\mathrm{quad},j}
  &=\sum_{\ell=0}^{24}\sum_{k=0}^{12}
       \frac{R_x^\ell\Delta^k}{\ell!k!}\rho_{j+\ell+2k},\\
 E_{\mathrm{Tay}}
  &=\mathcal M\left(\frac{R_x^{25}}{25!}
                  +e^{R_x}\frac{\Delta^{13}}{13!}\right).
                                                        \label{eq:Taylorerr}
\end{align}
Here \(R_x\le2,\ \Delta=5001/10^6\).
The midpoint order is at most \(49\) and the remainder order at most
\(51\); the supplied \(50\)- and \(52\)-order bounds cover them.
All Taylor segments have \(x\in[1000,1002]\) and
\(t\in[99999/10^6,21/200]\).

For explicit complete tails put \(\sigma=\sin(4d)\),
\(T_{\max}=1/5\), and
\[
 B=4c\pi m^2-9,\quad
 P_{\mathrm{gap}}=4\sigma\pi m^2-1002-2T_{\max}d,\quad
 r=((m+1)/m)^4e^{-c\pi(2m+1)}.
\]
Outward checks give \(B>190.5927,\ P_{\mathrm{gap}}>796.3701\),
\(r<0.000239\), and \(8c\pi m^2-T_{\max}>398.98\).
Define
\[
 S=\frac{(2\pi^2m^4+3\pi m^2)e^{-c\pi m^2}}{1-r},
 \quad Q_j=\sum_{k=0}^j{j\choose k}d^{j-k}k!/B^k .
\]
A normalized omitted-mode allowance is
\begin{equation}
                  E_{\mathrm{mode},j}=N_*\,2S Q_j/P_{\mathrm{gap}}.
                                                        \label{eq:modetail}
\end{equation}
To verify it, separate a shifted mode into its two monomials.
Their positive amplitudes have the form
\(g(u)=\exp(t(u^2-d^2)+bu-c\lambda_ne^{4u})\), \(b=9,5\).
For \(n\ge m\), the logarithm is concave,
\(g(u)\le g(0)e^{-Bu}\), and \(g'\le0\).
The negative phase derivative has magnitude
\(p(u)=4\sigma\lambda_ne^{4u}-x-2td\ge P_{\mathrm{gap}}>0\)
and \(p'>0\).
One integration by parts, including the derivative of
\((-\!d+\ii u)^j\), bounds the integral by
\[
 \frac{g(0)d^j}{p(0)}
 +\int_0^\infty
 \left\{\frac{(-g')(u+d)^j+jg(u+d)^{j-1}}p
                +\frac{g(u+d)^jp'}{p^2}\right\}\dd u .
\]
The integrand equals
\(-[g(u)(u+d)^j/p(u)]'+2jg(u)(u+d)^{j-1}/p(u)\).
Integrating it and using \(g\le g(0)e^{-Bu}\) gives
\(2g(0)Q_j/p(0)\).
The geometric mode sum, the two monomials, and the conjugate halves
with their factor \(1/2\) give \eqref{eq:modetail}.

For \(n<m\), the retained-source tail beyond \(U\) is bounded using
\[
 q_{n,j}(u)=(2\lambda_n^2+3\lambda_n)(u+d)^j
       e^{T_{\max}u^2+9u-c\lambda_ne^{4u}},\quad
 D_{n,j}=4c\lambda_ne^{4U}-9-2T_{\max}U-\frac j{U+d}.
\]
For \(j=0,1\), \(D_{n,j}>4121.56\), and the logarithm of \(q_{n,j}\)
is concave beyond \(U\). Its tangent gives
\begin{equation}
 E_{\mathrm{source},j}
      =N_*\sum_{n=1}^{11}\frac{q_{n,j}(U)}{D_{n,j}}.     \label{eq:spatialtail}
\end{equation}
The mode tail covers \(n\ge12\) on the whole shifted half-lines;
the retained-source tail covers \(n<12,u>2\). They are disjoint and
complete. The full normalized component allowance is
\[
 E_j=E_{\mathrm{quad},j}+E_{\mathrm{Tay}}+
                       E_{\mathrm{mode},j}+E_{\mathrm{source},j}.
\]
Parameter variation and arithmetic rounding are additionally contained
in the outward polynomial cell image itself.

\subsection{Recorded whole-cell output and interpretation}
\begin{center}\small
\begin{tabular}{lr}
\toprule
Complete normalized payment & Uniform upper bound\\
\midrule
Taylor remainder & \(2.083822530\,10^{-7}\)\\
Propagated quadrature radius & \(1.624269\,10^{-12}\)\\
Omitted-mode tail & \(1.784646\,10^{-12}\)\\
Retained-source tail & \(4.939296\,10^{-436}\)\\
Sum, rounded upward from displayed bounds & \(2.083856620\,10^{-7}\)\\
\bottomrule
\end{tabular}
\par\medskip
\begin{tabular}{lr}
\toprule
Closed \(x\)-cell & Lower bound \(H/G_*\)\\
\midrule
\([1001,1001.25]\) & \(0.01086203\)\\
\([1001.25,1001.5]\) & \(0.01153159\)\\
\([1001.5,1001.75]\) & \(0.00770777\)\\
\([1001.75,1002]\) & \(0.0004134538\)\\
\bottomrule
\end{tabular}
\end{center}
Each cell uses the entire closed time interval of \(\mathcal B\).
The lower bounds already pay every component error, parameter variation,
and arithmetic rounding. All seams and endpoints are included.
The separately recorded unrounded maximum added error is
less than \(2.083856619\,10^{-7}\); summing the outward-rounded rows
above justifies the slightly looser displayed total.
Every successful separating coordinate was \(H>0\); this result does
not need a nonzero derivative.

The 51 midpoint integrals were acquired with 160-bit Arb/Acb, normalized
absolute tolerance \(10^{-13}\), relative tolerance \(10^{-30}\),
evaluation limit \(200000\) and depth limit \(30\) per integral.
The positive-majorant integral used absolute tolerance \(10^{-22}\)
and relative tolerance \(10^{-25}\) before multiplication by \(N_*\),
with the same evaluation and depth limits.
The later wrapper consumed the same dyadic source cache at 192 and
256 bits, with four predeclared quarter-width cells, added-error cap
\(10^{-6}\), coordinate floor \(10^{-6}\), and a 300-second run cap.
There was no new midpoint quadrature, adaptive subdivision or scan.
Replaying two precisions is not independent source acquisition.
The own-data control
\[
 H_c=(x-1001.5)^2-2(t-1/10)
\]
has an event on the shared endpoints of two control cells.
The same strict first-jet consumer correctly refused event exclusion
on both. Theta error allowances were not assigned to this control.

Proposition~\ref{prop:box} is a direct value exclusion.
It removes a tiny time slice from the hypothetical event class; it
supplies no signed theta contour margin, no all-time strip, and no new
global bound on \(\Lambda\).

\section{A bounded continuation program}
The next numerical question should be fixed before any acquisition.
Choose one rational closed parameter rectangle remaining unresolved by
the existing coverage ledger; the box in Proposition~\ref{prop:box}
is already excluded. Fix a scale \(s_0>0\), normalization \(A_*>0\),
complete source domain, precision, cutoffs, allowed subdivisions, and
work budget. The primary candidate is a complete upper bound for
\(D_{s_0}/A_*^2\), compared with the strongest applicable direct
first-jet bound using the same data.

For finite cap integration with approximants \(P_0,P_1\) and complete
errors \(\varepsilon_0,\varepsilon_1\) on \(|y|\le L\), the product
payment includes
\[
 \int_{-L}^L G_s
 [s(2|P_1|\varepsilon_1+\varepsilon_1^2)
        +2|P_0|\varepsilon_0+\varepsilon_0^2]\dd y .
\]
If \(M_j\) bound the complete physical derivatives uniformly, a valid
Gaussian spatial tail is
\[
                  (sM_1^2+M_0^2)\erfc(L/(2\sqrt s)).
\]
At high height this crude allowance may dwarf the signal. Its size must
be assessed before acquiring a cap grid.
Alternatively use the signed source kernel \eqref{eq:Dkernel}.
Let \(K_D\) denote one-half of the bracket in \eqref{eq:Dkernel}.
If \(m_j\) and \(r_j\) are retained and omitted positive source moments,
\(|K_D|\le1+suv\) gives the complete pair-tail allowance
\[
                  2m_0r_0+r_0^2+s(2m_1r_1+r_1^2).
\]
Quadrature, rounding and parameter errors still require separate payment.
Source cutoffs in either variable must cover the entire omitted pair
region, including the overlap exactly once.

If \(U_D\) is the certified normalized deficit upper bound and
\(L_{\mathrm{direct}}\) is the strongest same-data lower bound for
\((h_0^2+s_0h_1^2)/A_*^2\), record
\[
 L_D=\max(0,-U_D),\qquad
 G_E=L_D-L_{\mathrm{direct}},\qquad \delta_D=-U_D.
\]
Only \(\delta_D>0\) certifies exclusion. Record acquisition cost,
incremental consumer cost, all error terms, and coverage as well as
the margin. The exact cap moments \eqref{eq:capmoments} belong in the
direct comparator; they must not be suppressed to manufacture a gain.
No such fixed window or completed deficit has yet been selected or
acquired.

Two secondary comparisons remain concrete:
a fixed nonnegative localized weight in \eqref{eq:localized} outside
the local positive region, and a proved nonzero-level arc with a complete
cap count and continuation contract. The former would contradict later
all-realness if negative; the latter would force a later event only
with at least four simple cap zeros under its full initial-event
accounting. Neither conclusion can be substituted for the other.

The new task sequence is: freeze a surviving window and source contract;
implement one finite consumer with own-data controls; execute one complete
theta comparison; review its margin or remaining error deficit; decide
whether another bounded attempt is justified. A failed complete estimate
is a legitimate scoped outcome. A changed scale, window or statistic
requires a new recorded contract. Root-interaction energy is deferred
until a complete or renormalized root system, exterior payments and
an independently estimable arithmetic statistic are specified.
Infinite-cap and higher-multiplicity contour work remain conditional
extensions, not consequences of the present finite tests.

\section*{Acknowledgments and research provenance}
Edward Baker proposed the expanding-contour, root-triangle splitting,
and heat flow terminology questions developed here.
This initial manuscript was drafted for him with substantial assistance
from OpenAI's gpt-6-astra through Codex, at recorded effort ultra, on
20 October 2026. Separate agents of the same model assisted derivation,
source checking and internal review. These statements disclose the
method of preparation; they do not constitute independent mathematical
verification.

This single source file consolidates route BW-09 Notes 01--06 and their
three internal reviews. The repository is
\url{https://github.com/ebbaker/shifted-zeta-positivity}.
The route is
\path{papers/investigations-3/13_contradiction_analysis/04_boundary_winding/threshold_routes/09_expanding_contour_later_collision/}.
The original research-progress register remains authoritative for
BW09-00--06. The separate heat-flow-contour and energy program/register
govern HFCE-00--07; their preparation creates no new quantitative result.

The exact polynomial implementation and receipt are
\path{numerics/check_expanding_contour_controls.py} and
\path{numerics/EXPANDING_CONTOUR_CONTROLS_20261020.json} within that route.
The bounded theta wrapper and receipt are
\path{numerics/check_bounded_first_jet_extension.py} and
\path{numerics/BOUNDED_THETA_EVENT_GATE_20261020.json}.
The inherited source cache and evaluator are in
\path{papers/investigations-3/11_arithmetic_collision_rigidity/numerics/high_height_consumer/}.
Source and review hashes are recorded in the route's
\path{numerics/BW09_SOURCE_VERSIONS_20261020.json} and
\path{numerics/BW09_REVIEW_RECORD_20261020.json}.
The numerical source acquisition's original manuscript manifest was
checked against commit
\path{ad4783b067548a10b8bfa65de8b4a0a47b1aa924};
later manuscript edits are not attributed to that acquisition.
These are provenance pointers, not dependencies needed to read the proofs.

\appendix
\section{Recorded dyadic source inputs}\label{app:cache}
% Exact dyadic input appendix for the BW09 manuscript.
% Prepared for Edward Baker, 20 October 2026, with substantial LLM assistance.
% Model: gpt-6-astra; reasoning effort: ultra.
% Requires the standard longtable package. Paste into the single TeX source;
% this file does not load external data or perform quadrature.

This appendix supplies all numerical inputs to the cached-moment
calculation. No decimal conversion or rounding of these inputs is used.
Set
\[
 d=\frac\pi8-\frac{4\log1000}{1000},\qquad
 t_c=\frac{21}{200},\qquad x_c=1000,\qquad
 \lambda_n=\pi n^2,\qquad v=u+id.
\]
The exact normalized finite-source midpoint derivatives are
\[
 I_j=1000^{9/4}\Re\int_0^2
 e^{t_cv^2+v+ix_cu}
 \sum_{n=1}^{11}
 (2\lambda_n^2e^{8v}-3\lambda_ne^{4v})
 e^{-\lambda_ne^{4v}}(iv)^j\,du.
\]
For each row of Table~\ref{tab:exact-source-input}, define
\[
 a_j=M_j2^{e_j},\qquad \rho_j=R_j2^{f_j}.
\]
The recorded outward quadrature assertion is
\[
                  I_j\in[a_j-\rho_j,a_j+\rho_j].
\]
The integers in the table are exact, including their signs. These are
160-bit source-acquisition enclosures; the subsequent 192-bit and
256-bit calculations reuse them without a new source acquisition.
The full list is given even though the bounded first-jet calculation
uses midpoint orders only through 49.

\begingroup
\scriptsize
\setlength{\tabcolsep}{3pt}
\begin{longtable}{r r r r r}
\caption{Exact dyadic midpoint and outward radius inputs.}
\label{tab:exact-source-input}\\
\hline
$j$ & $M_j$ & $e_j$ & $R_j$ & $f_j$\\
\hline
\endfirsthead
\multicolumn{5}{c}{\tablename\ \thetable\ (continued)}\\
\hline
$j$ & $M_j$ & $e_j$ & $R_j$ & $f_j$\\
\hline
\endhead
\hline
\multicolumn{5}{r}{Continued on next page}\\
\endfoot
\hline
\endlastfoot
0 & $-2263280269527565583007310275609432546331$ & $-135$ & $74212685$ & $-68$\\
1 & $187832141538870642634947848533804698138243$ & $-140$ & $110912977$ & $-69$\\
2 & $-68893550982494433405310433485739956939145$ & $-139$ & $613023105$ & $-71$\\
3 & $-398277603527930391830317925732144363527695$ & $-142$ & $251761319$ & $-70$\\
4 & $564082041797174262840121240162579156049131$ & $-141$ & $867249587$ & $-72$\\
5 & $-269447328700839526742101136483289565612025$ & $-141$ & $105910549$ & $-69$\\
6 & $-15396173261250504484832457636419584354098975$ & $-146$ & $551271437$ & $-72$\\
7 & $5895057015331406362244984468400152551040243$ & $-144$ & $95734937$ & $-69$\\
8 & $2109052848845200373147347007182981674093111$ & $-147$ & $934521689$ & $-72$\\
9 & $-63017518909490100866342084320844772236912371$ & $-147$ & $149189715$ & $-70$\\
10 & $11555774258321812674765839836042531484814991$ & $-145$ & $343063759$ & $-71$\\
11 & $11863258265468410131672668277813090377372269$ & $-145$ & $909851617$ & $-72$\\
12 & $-6124381681882501584912557284779999482074785$ & $-143$ & $693768569$ & $-72$\\
13 & $6175802700037614377584035864269358637720361$ & $-146$ & $514890109$ & $-71$\\
14 & $121448774021615353919706697833653958781231793$ & $-147$ & $127910091$ & $-70$\\
15 & $-110289578018384014547092804763233154419387595$ & $-147$ & $409912679$ & $-71$\\
16 & $-38862017812304035508391827356220322515980089$ & $-146$ & $927491085$ & $-72$\\
17 & $208118134565350579078023333997706081953035023$ & $-147$ & $504963671$ & $-72$\\
18 & $-56442462153430579955744620158491913724207569$ & $-147$ & $538391651$ & $-72$\\
19 & $-236010837405680444817311510368118666359425003$ & $-147$ & $131705037$ & $-70$\\
20 & $516565362574343969836451209057775785294912259$ & $-148$ & $578260685$ & $-72$\\
21 & $58760748504147526860067006723590484443002087$ & $-146$ & $39704377$ & $-68$\\
22 & $-28023062719467418591630986900448942452015062499$ & $-153$ & $692402693$ & $-72$\\
23 & $11543574767945077371673209606637445354166134167$ & $-153$ & $691564629$ & $-71$\\
24 & $28726620703596258282889747207215395599576049051$ & $-153$ & $510789769$ & $-71$\\
25 & $-18838482150507745117572142517625799607543454437$ & $-152$ & $907691401$ & $-72$\\
26 & $-9498833458096821636835979843025455116664856693$ & $-153$ & $536872561$ & $-72$\\
27 & $57900498313234543362401614019035041822054776505$ & $-153$ & $139491039$ & $-70$\\
28 & $-15941746186568209379734775010149494007433523521$ & $-152$ & $414118523$ & $-71$\\
29 & $-106242768322891097897297904820793144773794077481$ & $-154$ & $248465829$ & $-71$\\
30 & $41836570674964100729207930770647179613059614307$ & $-152$ & $156473249$ & $-70$\\
31 & $3421780527464540463252449497146518674286034297$ & $-152$ & $414325601$ & $-71$\\
32 & $-117450469900320997820779876468190710238103336031$ & $-153$ & $814687953$ & $-73$\\
33 & $20280380232682174927557859457428423396678891641$ & $-151$ & $571514109$ & $-72$\\
34 & $5941863038151970310378250792377969692120136563$ & $-149$ & $792149687$ & $-72$\\
35 & $-5678131967318140589304799714236950560101595787$ & $-148$ & $1053547513$ & $-73$\\
36 & $6056040011163460981476539139403998462927052635$ & $-152$ & $518106823$ & $-72$\\
37 & $117024517764362505680923289950802168881975707927$ & $-152$ & $329246989$ & $-71$\\
38 & $-48915829037666160527578506370698914573710499765$ & $-151$ & $634139765$ & $-72$\\
39 & $-163373710127647955370684856774196803625477365171$ & $-153$ & $185287249$ & $-70$\\
40 & $193492841869406967467357274736393108863731624355$ & $-152$ & $178607443$ & $-70$\\
41 & $-4850929722613871538595456957093956175949863165$ & $-149$ & $382481281$ & $-72$\\
42 & $-458309421882426708105541672934292862166982855949$ & $-153$ & $783305959$ & $-72$\\
43 & $28389340874355542127743407112124532591283914261$ & $-149$ & $185945327$ & $-72$\\
44 & $33214113921837982514899153774335369092695146723$ & $-150$ & $98831371$ & $-71$\\
45 & $-101272028942888213736150211654648027787665944961$ & $-150$ & $958122709$ & $-73$\\
46 & $65824704990555898454586870951875742050719048943$ & $-151$ & $738271163$ & $-72$\\
47 & $881727227277302018364380524409840958455155380585$ & $-153$ & $536951715$ & $-72$\\
48 & $-256016810219695535914779348969100370426413338365$ & $-151$ & $65699813$ & $-68$\\
49 & $-49659983180536780087676357917081185661747734835$ & $-150$ & $360719697$ & $-71$\\
50 & $208761838356981971367983994154359466972610209763$ & $-150$ & $173383791$ & $-70$\\
\end{longtable}
\endgroup

The positive normalized common derivative majorant is
\[
 \mathcal M=1000^{9/4}e^d\int_0^2
 \sum_{n=1}^{11}(2\lambda_n^2+3\lambda_n)
 e^{u^2/5+9u-c\lambda_ne^{4u}}
 [1+(u+d)^{52}]\,du,\qquad c=\cos(4d).
\]
Its recorded exact outward upper endpoint is
\[
 \mathcal M\le
 1024344365611329142069049401581335946080705817735\,2^{-123}
 <96328881210.
\]
Only an upper enclosure of this positive integral is required by the
Taylor remainder argument. This is the same exact dyadic majorant used
in both arithmetic replays; no display-rounding allowance is added or
removed.

The entries above were transcribed exactly from the pinned acquisition
cache. Its SHA-256 is
\begin{center}
\footnotesize\ttfamily
c5598f9dfcc9ba0601b85989703a1e945db54beaf91396258e9758a8a9f3604f
\end{center}
The source file is
\texttt{high\_height\_consumer/SETUP.json} in the accompanying numerical
artifacts. The transcribed row count, order, signed midpoint mantissas,
midpoint exponents, radius mantissas and radius exponents were all
checked for exact agreement with that cache. The cache's source and
historical manuscript pins are described in the provenance discussion.

This appendix removes the external cache as a source of missing
numerical inputs for the finite polynomial calculation. It does not,
by itself, prove the quadrature enclosure assertions: those assertions
are the results of the stated outward interval acquisition. Independent
verification can reacquire the explicitly defined integrals with
rigorous quadrature, then run the stated finite-error and cell algorithm.
A hash establishes data identity, not an integral inequality.


\begin{thebibliography}{9}
\bibitem{debruijn}
N.~G. de Bruijn, \emph{The roots of trigonometric integrals},
Duke Mathematical Journal \textbf{17} (1950), 197--226.
\url{https://doi.org/10.1215/S0012-7094-50-01720-0}.
\bibitem{newman}
C.~M. Newman, \emph{Fourier transforms with only real zeros},
Proceedings of the American Mathematical Society \textbf{61} (1976),
245--251.
\url{https://sites.math.northwestern.edu/~auffing/papers/Newman.pdf}.
\bibitem{rodgerstao}
B.~Rodgers and T.~Tao, \emph{The de Bruijn--Newman constant is non-negative},
Forum of Mathematics, Pi \textbf{8} (2020), e6.
\url{https://doi.org/10.1017/fmp.2020.6}.
\bibitem{nolen}
J.~Nolen, \emph{Partial Differential Equations and Diffusion Processes},
Stanford lecture notes, section on instantaneous heat propagation.
\url{https://math.stanford.edu/~ryzhik/STANFORD/STANF227-11/notes227-09.pdf}.
\bibitem{mccuan}
J.~McCuan, \emph{The Heat Equation: The Fundamental Solution and
Mean Value Property}, 19 April 2021.
\url{https://mccuan.math.gatech.edu/courses/6702/heateqn.pdf}.
\end{thebibliography}
\end{document}
